% ============================================================================
%  La firma digital y el viaje del mensaje firmado
%  Módulo 03 · Criptografía y redes — ChainNotes
%
%  Versión LaTeX del apunte web temas/firma-digital.html
%  Los ejemplos numéricos son los que produce assets/js/firma-digital.js
%  con la curva y^2 = x^3 + 3x + 6 mod 10007, n = 10039, G = (5,4270).
%
%  Compilar:  pdflatex firma-digital.tex   (dos pasadas, por el índice)
% ============================================================================

\documentclass[11pt,a4paper]{article}

% --- idioma y codificación -------------------------------------------------
\usepackage[T1]{fontenc}
\usepackage[utf8]{inputenc}
\usepackage[spanish,es-noshorthands,es-tabla]{babel}
\usepackage{lmodern}
\usepackage{microtype}

% --- página ----------------------------------------------------------------
\usepackage[a4paper,margin=2.6cm,headheight=14pt]{geometry}
\usepackage{fancyhdr}
\usepackage{enumitem}
\usepackage{booktabs}
\usepackage{tabularx}
\usepackage{array}
\usepackage{caption}
\usepackage{longtable}

% --- matemáticas -----------------------------------------------------------
\usepackage{amsmath,amssymb,amsthm}
\usepackage{mathtools}

% --- gráficas --------------------------------------------------------------
\usepackage{tikz}
\usetikzlibrary{positioning,arrows.meta,calc,fit,backgrounds,shapes.geometric,
                decorations.pathreplacing}

% --- enlaces y cajas (hyperref antes que tcolorbox) ------------------------
\usepackage[colorlinks=true,linkcolor=violetdeep,urlcolor=violetdeep,
            citecolor=violetdeep,pdftitle={La firma digital y la red P2P},
            pdfauthor={ChainNotes}]{hyperref}
\usepackage[most]{tcolorbox}

% ---------------------------------------------------------------------------
%  Paleta (heredada de assets/css/styles.css)
% ---------------------------------------------------------------------------
\definecolor{violetc}{HTML}{7A68F5}
\definecolor{violetdeep}{HTML}{5B46E0}
\definecolor{mintc}{HTML}{1FA88A}
\definecolor{coralc}{HTML}{E04E5C}
\definecolor{amberc}{HTML}{C98A1F}
\definecolor{skyc}{HTML}{1E88C7}
\definecolor{inkc}{HTML}{282344}

\hypersetup{linkcolor=violetdeep}

% ---------------------------------------------------------------------------
%  Estilo de párrafo y encabezados
% ---------------------------------------------------------------------------
\setlength{\parindent}{0pt}
\setlength{\parskip}{0.55em}
\captionsetup{font=small,labelfont={bf,color=violetdeep}}

\renewcommand{\topfraction}{0.9}
\renewcommand{\bottomfraction}{0.8}
\renewcommand{\textfraction}{0.07}
\renewcommand{\floatpagefraction}{0.75}
\setcounter{topnumber}{3}
\setcounter{totalnumber}{4}

\pagestyle{fancy}
\fancyhf{}
\fancyhead[L]{\small\color{inkc}\textsc{Firma digital y red P2P}}
\fancyhead[R]{\small\color{inkc}\thepage}
\renewcommand{\headrulewidth}{0.4pt}

% ---------------------------------------------------------------------------
%  Entornos de teoremas
% ---------------------------------------------------------------------------
\theoremstyle{definition}
\newtheorem{definicion}{Definición}
\theoremstyle{plain}
\newtheorem{proposicion}{Proposición}
\theoremstyle{remark}
\newtheorem*{observacion}{Observación}

% ---------------------------------------------------------------------------
%  Cajas de color
% ---------------------------------------------------------------------------
\newtcolorbox{nota}[1]{enhanced,breakable,colback=skyc!6,colframe=skyc!75!black,
  coltitle=white,fonttitle=\bfseries\small,title={#1},
  left=7pt,right=7pt,top=5pt,bottom=5pt,boxrule=0.7pt,arc=2.5pt,
  before upper={\setlength{\parskip}{0.5em}}}

\newtcolorbox{caso}[1]{enhanced,breakable,colback=amberc!8,colframe=amberc!85!black,
  coltitle=white,fonttitle=\bfseries\small,title={#1},
  left=7pt,right=7pt,top=5pt,bottom=5pt,boxrule=0.7pt,arc=2.5pt,
  before upper={\setlength{\parskip}{0.5em}}}

\newtcolorbox{aviso}[1]{enhanced,breakable,colback=coralc!7,colframe=coralc!85!black,
  coltitle=white,fonttitle=\bfseries\small,title={#1},
  left=7pt,right=7pt,top=5pt,bottom=5pt,boxrule=0.7pt,arc=2.5pt,
  before upper={\setlength{\parskip}{0.5em}}}

\newtcolorbox{clave}[1]{enhanced,breakable,colback=mintc!7,colframe=mintc!85!black,
  coltitle=white,fonttitle=\bfseries\small,title={#1},
  left=7pt,right=7pt,top=5pt,bottom=5pt,boxrule=0.7pt,arc=2.5pt,
  before upper={\setlength{\parskip}{0.5em}}}

\newtcolorbox{laboratorio}[1]{enhanced,breakable,colback=violetc!6,colframe=violetc!80!black,
  coltitle=white,fonttitle=\bfseries\small,title={#1},
  left=7pt,right=7pt,top=5pt,bottom=5pt,boxrule=0.7pt,arc=2.5pt,
  before upper={\setlength{\parskip}{0.5em}}}

\newtcolorbox{algoritmo}[1]{enhanced,breakable,colback=inkc!3,colframe=inkc!60,
  coltitle=white,fonttitle=\bfseries\small,title={#1},
  left=9pt,right=7pt,top=5pt,bottom=5pt,boxrule=0.6pt,arc=2.5pt}

% ---------------------------------------------------------------------------
%  Atajos de notación
% ---------------------------------------------------------------------------
\newcommand{\si}{\textcolor{mintc}{\checkmark}}
\newcommand{\no}{\textcolor{coralc}{$\times$}}
\newcommand{\mecanismo}[1]{\textbf{\textcolor{violetdeep}{#1}}}

\setlist[itemize]{label=\textcolor{violetc}{$\bullet$}}

\newlist{opciones}{enumerate}{1}
\setlist[opciones]{label=(\alph*),leftmargin=2.4em,itemsep=0.1em,topsep=0.2em}

\begin{document}

\makeatletter
\DeclareRobustCommand{\%}{\char`\%\relax}
\makeatother

\begin{titlepage}
\centering
\vspace*{2.5cm}

{\Large\color{violetdeep}\textsc{ChainNotes}\par}
{\small\color{inkc}Módulo 03 · Criptografía y redes\par}

\vspace{1.4cm}
{\Huge\bfseries La firma digital\par}
\vspace{0.4cm}
{\LARGE y el viaje del mensaje firmado\par}

\vspace{1.2cm}
\begin{minipage}{0.78\textwidth}
\centering\itshape\color{inkc}
Un nodo recibe una transacción de alguien a quien no conoce y tiene que decidir
solo, en milisegundos, si de verdad la autorizó quien dice. Eso lo resuelve la
firma digital. Pero una transacción que nadie recibe tampoco existe, así que la
otra mitad del problema es cómo ese mensaje firmado llega a toda la red. Las dos
mitades se encuentran en un punto: antes de reenviar nada, cada nodo verifica la
firma.
\end{minipage}

\vspace{1.6cm}
\begin{tikzpicture}[scale=1.0]
  \node[violetdeep,font=\small\bfseries] (d) at (-3.2,0) {$d$};
  \node[inkc,font=\scriptsize,below=1pt] at (d) {privada};
  \node[skyc,font=\small\bfseries] (Q) at (3.2,0) {$Q = d\cdot G$};
  \node[inkc,font=\scriptsize,below=1pt] at (Q) {pública};
  \draw[-{Stealth[length=5pt]},violetc,line width=1pt] (-2.7,0.25) -- (2.2,0.25);
  \node[violetc,font=\scriptsize,above] at (-0.25,0.25) {barato};
  \draw[-{Stealth[length=5pt]},coralc,line width=1pt,dashed] (2.2,-0.3) -- (-2.7,-0.3);
  \node[coralc,font=\scriptsize,below] at (-0.25,-0.3) {inviable};
\end{tikzpicture}

\vfill
{\small Apuntes de Blockchain · Proyecto académico\par}
{\small\color{inkc}\today\par}
\end{titlepage}

\tableofcontents
\clearpage

\begin{nota}{Cómo usar este documento}
Es la versión imprimible del apunte interactivo \texttt{temas/firma-digital.html}.
Los tres laboratorios de la página web aparecen aquí como ejemplos resueltos con
números fijos. La curva es deliberadamente pequeña ($p = 10\,007$) para que todas
las cuentas quepan en una línea y los ataques se puedan seguir a mano: la
inseguridad es intencional. Al final hay siete ejercicios con soluciones en el
apéndice~\ref{app:sol}.
\end{nota}

% ===========================================================================
\section{¿Quién autorizó esto?}\label{sec:quien}
% ===========================================================================

Una transacción llega a un nodo que no conoce a nadie y no puede preguntarle a
nadie. Tiene que decidir, él solo y en milisegundos, si ese mensaje lo autorizó
quien dice haberlo autorizado. No hay servidor central al que consultar, ni
sesión iniciada, ni soporte técnico.

Una contraseña no sirve aquí. Para comprobar una contraseña hay que conocerla, y
entonces quien la comprueba también puede usarla. En una red abierta eso
significa darle tu secreto a miles de desconocidos.

La firma digital resuelve exactamente esa asimetría: \textbf{firmar requiere un
secreto, verificar no requiere ninguno}.

\begin{clave}{Tres garantías distintas}
\textbf{Autenticidad:} el mensaje viene de quien dice.
\textbf{Integridad:} no fue alterado desde que se firmó.
\textbf{No repudio:} el firmante no puede negar después que fue él.

Las tres salen de la misma operación, pero conviene no mezclarlas: un hash solo
da integridad, y un MAC da autenticidad e integridad pero no repudio, porque
quien verifica también podría haberlo generado.
\end{clave}

% ===========================================================================
\section{Qué es un esquema de firma}\label{sec:esquema}
% ===========================================================================

\begin{definicion}[Esquema de firma]
Son tres algoritmos:
\[
  \mathrm{Gen}(1^\lambda) \to (d, Q), \qquad
  \mathrm{Firmar}(d, m) \to \sigma, \qquad
  \mathrm{Verificar}(Q, m, \sigma) \to \{0,1\}.
\]
\end{definicion}

\textbf{Correctitud:} una firma honesta siempre verifica,
$\mathrm{Verificar}(Q, m, \mathrm{Firmar}(d,m)) = 1$ para todo $m$.

\textbf{Seguridad (EUF-CMA).} El modelo estándar es la infalsificabilidad
existencial bajo ataque de mensaje elegido. Al adversario se le concede mucho:
conoce $Q$ y puede pedir la firma de todos los mensajes que quiera, elegidos por
él, incluso adaptándolos según lo que vaya viendo. Gana si produce una firma
válida de \emph{cualquier} mensaje que no haya pedido:
\[
  \Pr\bigl[\mathrm{Verificar}(Q, m^*, \sigma^*) = 1 \ \wedge\ m^* \notin \mathcal{Q}_{\text{firmadas}}\bigr]
  \;\le\; \mathrm{negl}(\lambda).
\]

No hace falta que el mensaje falsificado sea útil ni que signifique algo: basta
con uno. Esa exigencia tan estricta es deliberada, porque en un sistema real el
adversario elige qué le conviene falsificar.

\begin{observacion}
EUF-CMA no promete que la firma sea única, ni que no se pueda transformar en otra
firma válida del \emph{mismo} mensaje. Esa propiedad extra se llama
maleabilidad, y en Bitcoin costó cara: por eso existe SegWit, que sacó la firma
del cálculo del identificador de la transacción.
\end{observacion}

% ===========================================================================
\section{Se firma el hash, no el documento}\label{sec:hash}
% ===========================================================================

Las firmas sobre curvas operan con números menores que $n$, no con textos de
longitud arbitraria. Así que lo que se firma nunca es el mensaje: es $H(m)$. A
eso se le llama \emph{hash-then-sign}, y además de resolver el tamaño, hace la
firma más rápida y permite firmar sin ver el documento completo.

\begin{proposicion}[Firma con hash]
Sea $\Sigma$ un esquema EUF-CMA y $\Sigma_H$ el esquema que firma $H(m)$ en lugar
de $m$. Entonces $\Sigma_H$ es EUF-CMA si $\Sigma$ lo es \emph{y} $H$ es
resistente a colisión.
\end{proposicion}

\begin{aviso}{La clave privada nunca se compromete}
Si alguien encuentra $m_1 \ne m_2$ con $H(m_1) = H(m_2)$, puede pedirte que
firmes el inocuo y quedarse con una firma válida sobre el fraudulento. No se
rompió la curva, no se filtró $d$, no hubo fuerza bruta: se rompió la resistencia
a colisión del hash, y la firma cayó con ella.

Por eso Bitcoin firma sobre doble SHA-256, y por eso la migración de SHA-1 fue
urgente en cuanto las colisiones se volvieron prácticas en 2017.
\end{aviso}

% ===========================================================================
\section{Dónde vive la clave}\label{sec:curva}
% ===========================================================================

Una curva elíptica sobre un campo finito es el conjunto de puntos $(x,y)$ que
cumplen $y^2 = x^3 + ax + b \pmod p$, más un punto al infinito. Lo valioso es que
esos puntos forman un \textbf{grupo}.

Sumar un punto consigo mismo muchas veces es multiplicar por un escalar, y se
hace rápido duplicando: para calcular $1000\cdot G$ no hacen falta mil sumas sino
unas diez. La operación inversa —dado $G$ y $Q = d\cdot G$, hallar $d$— es el
\textbf{problema del logaritmo discreto en curva elíptica}, y es lo único que
impide que cualquiera deduzca tu clave privada a partir de tu dirección.

Bitcoin y Ethereum usan \textbf{secp256k1}, con un primo de 256 bits. El mejor
ataque conocido cuesta del orden de $2^{128}$ operaciones.

\begin{nota}{La curva de este documento}
\[
  y^2 = x^3 + 3x + 6 \pmod{10\,007}, \qquad n = 10\,039 \ \text{(primo)}, \qquad G = (5, 4270).
\]
Misma aritmética que secp256k1, cuatro cifras en vez de setenta y siete. Como $n$
es primo, cualquier punto distinto del infinito genera todo el grupo, así que no
hay claves débiles que esquivar.

Con $n = 10\,039$ tu clave privada se encuentra probando diez mil valores: cosa
de milisegundos. Eso es justamente lo que hace visibles los ataques de la
sección~\ref{sec:nonce}.
\end{nota}

% ===========================================================================
\section{ECDSA paso a paso}\label{sec:ecdsa}
% ===========================================================================

\subsection{Las ecuaciones}

\textbf{Firmar} con clave $d$ y nonce $k$:
\[
  R = k\cdot G, \qquad r = R_x \bmod n, \qquad
  e = H(m) \bmod n, \qquad
  s = k^{-1}(e + r\,d) \bmod n.
\]
La firma es el par $(r,s)$.

\textbf{Verificar} con la clave pública $Q$:
\[
  w = s^{-1}, \qquad P = (e\,w)\cdot G + (r\,w)\cdot Q,
  \qquad \text{válida} \iff P_x \bmod n = r.
\]

\begin{proposicion}[Por qué cierra la verificación]
Sustituyendo $Q = d\cdot G$ y sacando $G$ de factor común,
\[
  P = (e\,w)\cdot G + (r\,w)(d\cdot G) = \bigl(w(e + r\,d)\bigr)\cdot G.
\]
Como $w = s^{-1} = k\,(e + r\,d)^{-1}$, el paréntesis se cancela y
\[
  P = k\,(e+r\,d)^{-1}(e + r\,d)\cdot G = k\cdot G = R.
\]
El verificador reconstruye $R$ sin conocer $k$ ni $d$.
\end{proposicion}

\subsection{Ejemplo resuelto}

Con $d = 1234$, $k = 9871$ y el mensaje «pago 5 a Beto»:

\begin{center}
\small
\begin{tabular}{@{}ll@{}}
\toprule
Dato & Valor \\
\midrule
Clave pública $Q = d\cdot G$ & $(5354,\ 481)$ \\
$e = \mathrm{SHA256}(m) \bmod n$ & $5592$ \\
$R = k\cdot G$, de donde $r = R_x$ & $6126$ \\
$s = k^{-1}(e + r\,d) \bmod n$ & $7316$ \\
\midrule
Verificación & válida \\
\bottomrule
\end{tabular}
\end{center}

\begin{observacion}
Si cambias solo la clave privada y vuelves a firmar, $r$ no cambia: $r$ depende
únicamente de $k$. Si cambias el nonce, cambia todo. Esa dependencia es la que
explota el ataque de la sección siguiente.
\end{observacion}

\begin{aviso}{Dos secretos, no uno}
La clave privada es un secreto permanente. El nonce es un secreto de usar y
tirar, y es igual de crítico: filtrarlo o repetirlo entrega la clave privada
completa.
\end{aviso}

% ===========================================================================
\section{El nonce que entrega la clave}\label{sec:nonce}
% ===========================================================================

La $k$ de ECDSA tiene que ser secreta, impredecible y \textbf{distinta en cada
firma}. Las tres condiciones son obligatorias, y la tercera es la que más se
incumple, porque incumplirla no produce ningún error visible: las firmas siguen
siendo válidas.

\subsection{El despeje}

Dos firmas con la misma $k$ sobre mensajes distintos comparten la misma $r$:
\[
  s_1 = k^{-1}(e_1 + r\,d), \qquad s_2 = k^{-1}(e_2 + r\,d).
\]
Restando, el término con $d$ desaparece:
\[
  s_1 - s_2 = k^{-1}(e_1 - e_2)
  \quad\Longrightarrow\quad
  k = \frac{e_1 - e_2}{s_1 - s_2} \bmod n,
\]
y con $k$ en la mano sale la clave privada:
\[
  d = \frac{s_1 k - e_1}{r} \bmod n.
\]

\begin{laboratorio}{Ejemplo resuelto}
Con $d = 1234$ y $k = 4242$ reutilizado en «pago 5 a Beto» y «pago 7 a Caro»:

\begin{center}
\small
\begin{tabular}{@{}lrr@{}}
\toprule
 & Firma 1 & Firma 2 \\
\midrule
$r$ & 6153 & 6153 \\
$s$ & 7542 & 5657 \\
\bottomrule
\end{tabular}
\end{center}

Las dos $r$ coinciden en $6153$, lo que delata públicamente el reuso. El despeje da
$k = 4242$ y $d = 1234$: la clave privada real.
\end{laboratorio}

\begin{observacion}
No interviene el tamaño de la curva. Con secp256k1 el ataque es igual de
instantáneo: 256 bits no protegen de un error de implementación. No se resuelve
ningún logaritmo discreto, son dos ecuaciones lineales en $\mathbb{Z}_n$ con dos
incógnitas.
\end{observacion}

\subsection{No es teórico}

En 2010 Sony firmaba el software de la PlayStation 3 con una $k$ constante, y el
grupo fail0verflow extrajo la clave maestra de la consola. En 2013 un generador
de azar defectuoso en Android repitió nonces en carteras de Bitcoin y se vaciaron
monederos reales. En los dos casos la curva era sólida y la implementación de la
firma correcta; lo que falló fue el número aleatorio.

\begin{clave}{La solución: no usar azar}
El \textbf{RFC 6979} deriva $k$ de forma determinista a partir del mensaje y de
la clave privada, con un HMAC. Así $k$ sigue siendo impredecible para quien no
tenga la clave, pero deja de depender del generador aleatorio del dispositivo.
Es lo que usan hoy las carteras serias, y la razón por la que dos firmas del
mismo mensaje con la misma clave salen idénticas.
\end{clave}

% ===========================================================================
\section{Schnorr y la multifirma}\label{sec:schnorr}
% ===========================================================================

Bitcoin añadió Schnorr en 2021 con Taproot. El motivo no fue la seguridad sino
una propiedad algebraica que ECDSA no tiene.

\textbf{Firmar:}
\[
  R = k\cdot G, \qquad e = H(R \parallel Q \parallel m), \qquad
  s = k + e\,d \bmod n.
\]
\textbf{Verificar:} $s\cdot G = R + e\cdot Q$.

Comparada con la de ECDSA llama la atención lo que le falta: no hay ningún
$k^{-1}$. Es \textbf{lineal} en la clave y en el nonce, y de ahí sale todo lo
demás.

\subsection{Agregación}

Si dos personas firman el mismo mensaje, sus firmas se suman:
\[
  Q = Q_1 + Q_2, \qquad R = R_1 + R_2, \qquad s = s_1 + s_2,
\]
y el resultado verifica contra la suma de las claves públicas, porque
\[
  s_1 + s_2 = (k_1 + k_2) + e(d_1 + d_2)
\]
tiene exactamente la forma $s = k + e\,d$ con $k$ y $d$ agregados. Para la red
eso es indistinguible de una firma de un solo firmante: mismo tamaño, mismo costo
de verificación. Un monedero de tres de tres deja de anunciarse como tal.

\begin{observacion}
Con ECDSA no funciona. Cada sumando lleva su propio $k^{-1}$ y no hay forma de
factorizar uno común, porque $k_1^{-1} + k_2^{-1} \ne (k_1+k_2)^{-1}$. Además
cada firma tiene su propia $r$, mientras que una firma agregada necesitaría una
sola. Esa es la diferencia completa entre los dos esquemas.
\end{observacion}

\subsection{El ataque de clave deshonesta}

\begin{aviso}{Sumar claves públicas tal cual está roto}
Quien publica su clave al último puede elegirla en función de las ajenas. Si Beto
ve la clave $Q_1$ de Alicia y publica
\[
  Q_2 = d\cdot G - Q_1,
\]
entonces la clave agregada queda
\[
  Q = Q_1 + Q_2 = d\cdot G,
\]
que es exactamente su propia clave, sin rastro de Alicia. Beto firma solo y la
firma verifica contra la clave de los dos. Alicia nunca participó y no se entera.
\end{aviso}

Por eso lo que se usa de verdad es \textbf{MuSig2}, que obliga a cada firmante a
comprometerse primero con su nonce y pondera cada clave con un coeficiente que
depende de todas las demás, de modo que nadie puede elegir la suya en función de
las ajenas.

% ===========================================================================
\section{De la firma a la red}\label{sec:puente}
% ===========================================================================

Hasta aquí el mensaje está firmado, pero sigue en tu computadora. Una transacción
que nadie recibe no existe. La otra mitad del problema es cómo ese mensaje
firmado llega a miles de nodos que no se conocen entre sí.

\begin{clave}{Dónde se encuentran las dos mitades}
Antes de reenviar nada, \textbf{cada nodo verifica la firma}. Si no cuadra, tira
el mensaje y no lo propaga.

Sin firmas, el gossip sería un canal de inundación: cualquiera podría generar
transacciones falsas a nombre ajeno y la red las repartiría obedientemente. La
firma convierte cada nodo en un filtro que cuesta cero confianza, porque
verificar no requiere conocer a nadie ni consultar a nadie. Es lo que permite que
la red no tenga portero.
\end{clave}

\subsection{Cómo se propaga}

Cada nodo mantiene conexiones con un puñado de pares y les reenvía lo que recibe
y valida. En pocos saltos el mensaje alcanza a toda la red, sin que exista un
directorio central ni un orden de reparto.

Las transacciones válidas que aún no están en un bloque esperan en la
\textbf{mempool} de cada nodo. No es una cola global: cada nodo tiene la suya, y
pueden diferir.

\subsection{Qué ocurre si aparecen dos bloques}

Dos mineros pueden resolver el acertijo casi a la vez. Entonces parte de la red
ve un bloque y parte ve el otro: es una \textbf{bifurcación temporal}. Se resuelve
sola en cuanto aparece el siguiente bloque, porque todos adoptan la cadena con
más trabajo acumulado y la otra rama queda huérfana.

\begin{observacion}
Por eso la finalidad en una red de prueba de trabajo es \emph{probabilística}: no
hay un instante en que el bloque se declare definitivo, solo una probabilidad de
reversión que cae con cada confirmación.
\end{observacion}

% ===========================================================================
\section{Síntesis}\label{sec:sintesis}
% ===========================================================================

\begin{center}
\small
\begin{tabularx}{\textwidth}{@{}lX@{}}
\toprule
Pregunta & La responde \\
\midrule
¿Quién autorizó esto?        & La firma: solo quien tiene $d$ produce $(r,s)$; cualquiera lo comprueba con $Q$. \\
¿Lo alteraron?               & La firma: se firma $H(m)$, así que cambiar un carácter rompe la verificación. \\
¿Puede negarlo después?      & La firma: verificar no requiere el secreto, así que nadie más pudo firmarlo. \\
¿Cómo se entera la red?      & El gossip: cada nodo reenvía a sus pares. \\
¿Por qué no se inunda de basura? & Las dos: cada nodo verifica \emph{antes} de reenviar. \\
¿Y si llegan dos bloques a la vez? & El consenso: bifurcación temporal y cadena con más trabajo. \\
\bottomrule
\end{tabularx}
\end{center}

\begin{enumerate}[leftmargin=1.5em]
  \item Firmar necesita un secreto; verificar no necesita ninguno. Esa asimetría
        permite una red abierta sin portero.
  \item Nunca se firma el documento, sino su hash. La seguridad de la firma queda
        amarrada a la resistencia a colisión del hash.
  \item ECDSA necesita un nonce secreto, impredecible y distinto cada vez.
        Repetirlo entrega la clave privada con álgebra elemental.
  \item Schnorr es lineal, así que las firmas se suman. La agregación ingenua de
        claves es atacable; por eso existe MuSig2.
  \item Una transacción firmada que nadie recibe no existe.
\end{enumerate}

% ===========================================================================
\section{Ejercicios}\label{sec:ej}
% ===========================================================================

\begin{enumerate}[leftmargin=1.8em,itemsep=0.8em]

\item\label{e:verif} \textbf{Por qué cierra la verificación.}
Partiendo de $s = k^{-1}(e + r\,d)$, demuestra que el punto que calcula el
verificador es exactamente $R = k\cdot G$. Explica en una línea por qué eso lo
convence de que quien firmó conocía $d$.

\item\label{e:nonce} \textbf{El nonce repetido.}
Deduce la fórmula que recupera $k$ a partir de dos firmas que lo comparten, y
luego la que recupera $d$. Explica por qué el tamaño de la curva no influye en la
dificultad del ataque.

\item\label{e:agregar} \textbf{Por qué ECDSA no se puede agregar.}
Escribe $s_1 + s_2$ en ECDSA y explica qué término impide que el resultado tenga
la forma de una firma ECDSA válida.

\item\label{e:rogue} \textbf{La clave deshonesta.}
Verifica que si Beto publica $Q_2 = d\cdot G - Q_1$, la clave agregada es
exactamente $d\cdot G$. ¿Qué tendría que impedir un protocolo para evitarlo?

\item\label{e:mac} \textbf{Firma contra MAC.}
Un MAC da autenticidad e integridad y es mucho más barato que una firma. Explica
por qué Bitcoin no puede usar MAC en lugar de firmas.

\item\label{e:colision} \textbf{Colisión y falsificación.}
Describe un ataque concreto contra un esquema \emph{hash-then-sign} cuyo hash
admite colisiones, sin tocar en ningún momento la clave privada.

\item\label{e:gossip} \textbf{Sin firmas no hay gossip.}
Explica qué le pasaría a una red P2P que propagara transacciones sin verificar
firmas, y por qué el problema no se arregla poniendo un nodo de confianza que las
revise.

\end{enumerate}

% ===========================================================================
\appendix
\section{Soluciones}\label{app:sol}
% ===========================================================================

\subsection*{Ejercicio \ref{e:verif}}
Sustituyendo $Q = d\cdot G$ y factorizando,
$P = (e\,w)G + (r\,w)(dG) = \bigl(w(e+rd)\bigr)G$. Como
$w = s^{-1} = k(e+rd)^{-1}$, queda $P = kG = R$, y por eso $P_x \bmod n = r$.

Lo convence porque la igualdad solo puede cumplirse si quien construyó $s$
conocía el $d$ que satisface $Q = dG$; producir una $s$ que cuadre sin saber $d$
equivale a resolver el logaritmo discreto.

\subsection*{Ejercicio \ref{e:nonce}}
De $s_1 = k^{-1}(e_1 + rd)$ y $s_2 = k^{-1}(e_2 + rd)$, restando se cancela $d$:
$s_1 - s_2 = k^{-1}(e_1-e_2)$, de donde $k = (e_1-e_2)(s_1-s_2)^{-1}$. Despejando
de la primera, $d = (s_1 k - e_1)\,r^{-1}$.

El tamaño de la curva no influye porque no se resuelve ningún logaritmo discreto:
son dos ecuaciones lineales en $\mathbb{Z}_n$, y resolverlas cuesta un par de
inversos modulares. La seguridad de la curva protege contra el problema difícil,
no contra un error de implementación que regala la información.

\subsection*{Ejercicio \ref{e:agregar}}
En Schnorr, $s_1+s_2 = (k_1+k_2) + e(d_1+d_2)$, que tiene la forma de una firma
con clave y nonce agregados. En ECDSA,
\[
  s_1 + s_2 = k_1^{-1}(e + r_1 d_1) + k_2^{-1}(e + r_2 d_2),
\]
y el estorbo es el \textbf{inverso del nonce}: no hay forma de factorizar un
$k^{-1}$ común porque $k_1^{-1} + k_2^{-1} \ne (k_1+k_2)^{-1}$. Además cada firma
tiene su propia $r$, mientras que una agregada necesitaría una sola.

\subsection*{Ejercicio \ref{e:rogue}}
Por construcción $Q_1 + Q_2 = Q_1 + (dG - Q_1) = dG$. La suma de puntos es
conmutativa y $Q_1$ se cancela con su inverso. Beto conoce $d$, así que firma
solo y la verificación $sG = R + eQ$ se cumple.

Un protocolo lo evita obligando a cada firmante a \textbf{comprometerse con su
clave y su nonce antes de ver los ajenos}, y ponderando cada clave con un
coeficiente derivado del conjunto completo, de modo que nadie pueda elegir la
suya en función de las demás. Eso es MuSig2.

\subsection*{Ejercicio \ref{e:mac}}
Un MAC usa una clave \emph{compartida}: quien verifica también puede generar. En
una red abierta eso significa que cada nodo que valide tu transacción podría
fabricar otras a tu nombre, y además no habría no repudio, porque no se podría
distinguir quién de los dos la creó. La firma digital separa generar de
verificar, y esa separación es justo lo que la red necesita.

\subsection*{Ejercicio \ref{e:colision}}
El atacante construye $m_1$ inocuo y $m_2$ fraudulento con $H(m_1) = H(m_2)$. Te
pide que firmes $m_1$; tú lo lees, te parece inofensivo y lo firmas, produciendo
$\sigma$. Como la firma se calculó sobre $H(m_1) = H(m_2)$, el par $(m_2,\sigma)$
verifica igual de bien. La clave privada nunca salió de tu poder.

\subsection*{Ejercicio \ref{e:gossip}}
Sin verificación, cualquiera podría inundar la red con transacciones falsas a
nombre ajeno, y los nodos las repartirían obedientemente hasta agotar ancho de
banda y memoria. Es una denegación de servicio gratuita para el atacante.

Un nodo de confianza que las revise no arregla el problema: reintroduce un punto
central que puede censurar, equivocarse o caerse, y obliga a todos los demás a
confiar en él. La gracia de la firma es que \emph{cada} nodo puede filtrar por su
cuenta sin confiar en nadie.

% ===========================================================================
\section{Referencias}
% ===========================================================================

\begin{itemize}[leftmargin=1.4em]
  \item Goldwasser, S., Micali, S. y Rivest, R. \emph{A Digital Signature Scheme
        Secure Against Adaptive Chosen-Message Attacks}, SIAM J. Comput., 1988.
        Define EUF-CMA.
  \item Johnson, D., Menezes, A. y Vanstone, S. \emph{The Elliptic Curve Digital
        Signature Algorithm (ECDSA)}, IJIS, 2001.
  \item Schnorr, C.~P. \emph{Efficient Signature Generation by Smart Cards},
        Journal of Cryptology, 1991.
  \item NIST, \emph{FIPS 186-5: Digital Signature Standard}, 2023.
  \item Pornin, T. \emph{RFC 6979: Deterministic Usage of DSA and ECDSA}, 2013.
  \item BIP-340 (Schnorr para secp256k1) y BIP-341 (Taproot), activados en
        noviembre de 2021.
  \item Nick, J., Ruffing, T. y Seurin, Y. \emph{MuSig2: Simple Two-Round Schnorr
        Multi-Signatures}, CRYPTO 2021.
  \item fail0verflow, \emph{Console Hacking 2010}, 27C3. La clave de la PS3.
  \item Nakamoto, S. \emph{Bitcoin: A Peer-to-Peer Electronic Cash System}, 2008,
        secciones 5 y 8.
\end{itemize}

\end{document}
